\begin{corollary}[4.1.4]
\label{IV.2.4.1.4-fr}
Sean $X$ un preesquema localmente de tipo finito sobre un cuerpo $k$, $K$ una
extensión de $k$; entonces $\dim(X\otimes_k K)=\dim(X)$.
\end{corollary}

\begin{proof}
Es evidente que basta restringirse al caso en que $X$ sea de tipo finito sobre
$k$; entonces el morfismo $u:X\otimes_k K\to X$ es fielmente plano (2.2.13,
(i)), por tanto, si $U$ es un abierto en todas partes denso en $X$,
$u^{-1}(U)=U\otimes_k K$ es denso en $X\otimes_k K$ (2.3.10); si
$g:U\to\mathbf V_k^n$ es finito y sobreyectivo, lo mismo ocurre con
$g_{(K)}:U\otimes_k K\to\mathbf V_K^n$ (\textbf I, 3.5.2 y \textbf{II}, 6.1.5), de donde se
 sigue el corolario.
\end{proof}

\begin{corollary}[4.1.5]
\label{IV.2.4.1.5-fr}
Sean $X$ e $Y$ dos preesquemas localmente de tipo finito sobre un cuerpo $k$;
entonces $\dim(X\times_k Y)=\dim(X)+\dim(Y)$.
\end{corollary}

\begin{proof}
Basta demostrar que si $U$ (resp. $V$) es una vecindad afín de un punto $x$
(resp. $y$) de $X$ (resp. $Y$) tal que
$\operatorname{deg.tr}_k \boldsymbol{k}(x)=m=\dim U$ y
$\operatorname{deg.tr}_k \boldsymbol{k}(y)=n=\dim V$, entonces se tiene
$\dim(U\times_k V)=m+n$, en otras palabras, en virtud de (4.1.2), se puede
suponer que existen $k$-morfismos finitos sobreyectivos
$f:X\to\mathbf V_k^m$, $g:Y\to\mathbf V_k^n$; entonces
$f\times g:X\times_k Y\to\mathbf V_k^{m+n}$ es finito y sobreyectivo
(\textbf I, 3.5.2 y \textbf{II}, 6.1.5), de donde se sigue el corolario.
\end{proof}
